% Methods — ReCAHS paper

\section{Methods}
\label{sec:methods}

\subsection{Models, data and protocol}
\label{sec:methods-setup}
Our primary model is PatchTST \citep{nie2023patchtst} with
$d_{\text{model}}{=}128$, $d_{\text{ff}}{=}256$, $3$ encoder layers of
$8$ heads ($H{=}24$ heads total), input length $336$, trained with the
Time-Series-Library recipe (Adam, lr $10^{-4}$, batch $32$, $10$ epochs,
early stopping patience $3$, lr halved per epoch). We evaluate on ETTh1,
ETTh2, ETTm1 and Weather with prediction horizons
$\{96,192,336,720\}$ and $5$ training seeds per cell
($80$ models). As a generalization test we repeat the core pipeline on
iTransformer \citep{liu2024itransformer} at horizon $96$ with the same
capacity budget, noting that its heads attend across \emph{variates}
rather than time. All experiments run in a single pinned environment
(library commit, dataset and checkpoint hashes recorded); validation and
test loaders are unshuffled so window indices align with regime labels.

\subsection{Regime labels}
\label{sec:methods-regimes}
Each length-$336$ input window of the target variable is decomposed with
robust STL \citep{cleveland1990stl} (period $=$ one day at the dataset's
sampling rate). With component variances normalized to sum to one, the
window's regime is the argmax of
$(s_{\text{trend}}, s_{\text{seasonal}}, s_{\text{residual}})$, and the
\emph{confidence margin} is the gap between the two largest scores.
Labels are computed independently for train, validation and test splits;
test regime mixes range from strongly trend-skewed (ETTh1: $93.6\%$
trend) to strongly seasonal (Weather: $\geq 84\%$ seasonal).

\subsection{Head masking and selection criteria}
\label{sec:methods-selection}
A head $(\ell,h)$ is deactivated by zeroing its attention output before
the layer's output projection; this functional masking is numerically
equivalent to structurally removing the head's parameter slices. All
methods select masks on \emph{validation} data only:
\begin{itemize}
  \item \textbf{Independent importance (static).} Leave-one-out scoring
  \citep{michel2019sixteen}: importance of $(\ell,h)$ is the increase in
  validation MSE when it alone is masked; the lowest-scoring heads are
  pruned. Removal order $=$ ascending importance.
  \item \textbf{Joint greedy (pooled).} Greedy backward elimination
  \citep{li2021differentiable}: starting from all heads, repeatedly
  remove the head whose removal least harms validation MSE \emph{given
  the removals so far}, on a fixed $256$-window validation subsample.
  \item \textbf{Joint greedy (per-regime).} The same procedure restricted
  to validation windows of one regime, yielding three specialist masks;
  the dynamic policy applies the mask matching each window's detected
  regime.
  \item \textbf{Baselines.} Random head subsets ($3$ draws per seed) and
  magnitude ranking by the $\ell_2$ norm of each head's projection
  slices.
\end{itemize}
Removal orders are recorded down to $6$ remaining heads, so a mask at any
sparsity in $\{12.5,25,37.5,50,62.5,75\}\%$ is a prefix; all methods are
compared at matched sparsity. Every evaluation stores per-window test and
validation losses, enabling all subsequent analyses --- routing, oracles,
ensembling, statistics --- without re-running models.

\subsection{Routing, oracle and ensembles}
\label{sec:methods-routing}
At $25\%$ sparsity the candidate set is
$\mathcal{M}=\{$trend, seasonal, residual, pooled-greedy$\}$. The
\textbf{oracle} routes each test window to its loss-minimizing mask
(uses test targets; an upper bound, not a method). Deployable routers,
all trained on validation only: (i)~hard STL routing;
(ii)~a $12$-feature gradient-boosted \textbf{regime detector} replacing
STL at inference; (iii)~\textbf{learned oracle-routers} --- argmin
classification, per-mask loss regression, and centered (regret)
regression with margin weighting; (iv)~\textbf{delayed-feedback}
routing on observed regret of windows old enough for their targets to be
available, and its sliding-window bandit variants;
(v)~\textbf{regret-gated} deviation from the pooled mask. Failure causes
are separated by three diagnostics: in-sample fits (capacity),
a learning curve over $2.2$k--$27$k router-training windows using
additionally computed \emph{train-split} mask losses (data), and
fit-on-half-of-test plus routing from mean-pooled encoder
representations (shift / feature poverty).

\textbf{Ensembles} replace selection by combination: uniform averaging of
the four masked predictions, the top-$2$ masks by validation MSE,
weights $\propto$ inverse validation MSE, and \emph{soft regime routing}
in which the detected regime's mask receives weight $\alpha$ (others
share $1-\alpha$), with $\alpha{=}1$ recovering hard routing. A
\textbf{diversity ablation} rebuilds the pool with greedy masks from
random validation subsamples (no regime information) under identical
settings.

\subsection{Retraining}
\label{sec:methods-scratch}
To test whether pruning cost is intrinsic, the $25\%$ masks are imposed
\emph{from initialization} and the model trained under the standard
recipe (static mask; and a regime-conditional variant in which each
training window applies its regime's mask). This complements the earlier
finding that fine-tuning an already-converged checkpoint under the mask
fails.

\subsection{Statistical analysis}
\label{sec:methods-stats}
Within each cell, methods are compared by seed-paired differences
(Wilcoxon) and across the $16$ cells by two-sided sign tests. Cell-level
uncertainty uses a hierarchical circular moving-block bootstrap
(resampling seeds and $96$-window blocks, $10^4$ replicates, $95\%$
CIs). Regime dependence of head importance is tested by circular
rotations of the regime-label sequence, which preserve temporal
autocorrelation. Reported oracle-gain \emph{recovery} of a method $m$ is
$(L_{\text{pooled}}-L_m)/(L_{\text{pooled}}-L_{\text{oracle}})$.
